\section{Saddle Point System}
Saddle point system arises frequently in scientific and engineering applications, including 
constrained optimization problem, image reconstruction and optimal control.
Saddle point system has been extensively and systematically studied, notably
by~\cite{Benzi_Golub_Liesen_2005}.
In our case, the saddle point system arises naturally from the mixed finite 
element discretization.

In this section,
we give a brief overview of the properties of a saddle point 
systems and present an iteration method tailored to our problem.
We express the abstract saddle point system as follows:
\begin{equation}\label{eq:saddle}
		  \begin{bmatrix}
					 \bold{A} & \bold{B} \\
					 \bold{B}^T &\bold{-C}
		  \end{bmatrix}\begin{bmatrix}
					 \bold{x}\\
					 \bold{y}
		  \end{bmatrix} = \begin{bmatrix}
					 \bold{F}\\
					 \bold{G}
		  \end{bmatrix},
\end{equation}
where we assume the following properties:
\begin{itemize}
		  \item $\bold{A}$ is symmetric: $\bold{A} = \bold{A}^T$;
		  \item $\bold{A}$ is nonsingular: $\bold{A}^{-1}$ exists;
		  \item $\bold{C}$ is symmetric and positive semidefinite;
		  \item $\bold{B}$ is rank deficient.
\end{itemize}

\subsection{Reduced system}
We now present a detailed discussion of the saddle point system (SPS) arising from the mixed finite element method (MFEM)
including the treatment of boundary conditions.

As a first step, we regroup the degrees of freedom (dofs) ($i\in I$) into three distinct sets: the interior dofs $I_c$, 
the boundary dofs assigned with essential boundary condition $I_{\Gamma_u}$ and the boundary dofs assigned with 
natural boundary condition $I_{\Gamma_i}$. We have $I_c \cap I_{\Gamma_u} = I_c \cap I_{\Gamma_i} 
= I_{\Gamma_u} \cap I_{\Gamma_i} = \O$ and $I_c \cup I_{\Gamma_i} \cup I_{\Gamma_u} = I$.
We then reorganize the $\mathbf{A}$ matrix as
\begin{equation}\label{eq:full_system}
		  \mathbf{A} = 
		  \begin{bmatrix}
					 \mathbf{A}_K  & \mathbf{A}_M \\
					 \mathbf{A}_M^T & \mathbf{M}  \\
		  \end{bmatrix},
\end{equation}
where $\mathbf{A}_K = \{a_{k_1,k_2}\}$ with $k_1,k_2 \in I_c+I_{\Gamma_i}$, consists of all the ``free" dofs.
The matrix $\mathbf{A}_{M} = \{a_{m_1,m_2}\}$ with $m_1\in I_c+I_{\Gamma_i},\, m_2 \in I_{\Gamma_u}$, represents the 
interaction between ``free" and prescribed dofs.
Finally, $\mathbf{M} = \{m_{i,j}\},\, i,j \in I_{\Gamma_u}$ accounts for the contributions from
the fixed dofs associated with essential boundary conditions.

We perform the same treatment for matrix $\mathbf{B}$ and right hand
side vector $\mathbf{f}$, by regrouping
dofs according to its iteration with boundary conditions.
\begin{equation}
		  \mathbf{B} = \begin{bmatrix}
					 \mathbf{B}_K \\
					 \mathbf{B}_M
		  \end{bmatrix} ,\text{and}\quad \mathbf{F} = \begin{bmatrix}
					 \mathbf{F}_K\\
					 \mathbf{F}_M
		  \end{bmatrix}
\end{equation}
We rewrite the linear system with the regrouped dofs and form a reduced system
as
\begin{equation}
		  \begin{bmatrix}
					 \mathbf{A}_K & \mathbf{B}_K\\
					 \mathbf{B}^T_K & -\mathbf{C}
		  \end{bmatrix} \begin{bmatrix}
					 \mathbf{x} \\
					 \mathbf{y}
		  \end{bmatrix} = \begin{bmatrix}
					 \mathbf{F}_K - \mathbf{A}_M \mathbf{g} - \mathbf{g}_N\\
					 \mathbf{G} - \mathbf{B}_M^T \mathbf{g}
		  \end{bmatrix},
\end{equation}
where $\mathbf{g}$ denotes the values prescribed for essential 
boundary, and $\mathbf{g}_N$ is the vector generated by the natural boundary
condition part.
The reformed reduced system is still a saddle point
syste.


\subsection{Linear system for coupled Darcy-Stokes equations}
		  %The linear system formed directly from the equations~\eqref{eq:weak_stokes_mass}, 
		  %~\eqref{eq:weak_compaction_mass}, ~\eqref{eq:weak_darcy_mass} 
		  %and~\eqref{eq:weak_continuous_mass} is
The linear system formed directly from the equations~\eqref{eq:weak_stokes_mass}--\eqref{eq:weak_continuous_mass} is shown as
\begin{equation}
\begin{bmatrix}
\bold{A}_s   & -\bold{B}_s & \bold{0}     &  \bold{0}   \\
\bold{B}_s^T &  \bold{C}_s & \bold{0}     &  \bold{K}   \\
\bold{0}     &  \bold{0}   & \bold{A}_d   & -\bold{B}_d \\
\bold{0}     &  \bold{K}   & \bold{B}_d^T &  \bold{C}_d
\end{bmatrix} \begin{bmatrix}
\bold{x}_s\\
\bold{y}_s\\
\bold{x}_d\\
\bold{y}_d
\end{bmatrix}= 
\begin{bmatrix}
\bold{f}_s\\
\bold{g}_s\\
\bold{f}_d\\
\bold{g}_d
\end{bmatrix},
\end{equation}
wherein the solution represents the degrees of freedom of primary variables with respect to the bases of the finite element spaces.
(Insert details of calculation of each blocked matrix).
We form a double saddle point system by symmetrically permutate original linear system,
\begin{equation}
\begin{bmatrix}
\bold{A}_s   & \bold{0}     & -\bold{B}_s &  \bold{0}   \\
\bold{0}     & \bold{A}_d   &  \bold{0}   & -\bold{B}_d \\
\bold{B}_s^T & \bold{0}     &  \bold{C}_s &  \bold{K}   \\
\bold{0}     & \bold{B}_d^T &  \bold{K}   &  \bold{C}_d
\end{bmatrix} \begin{bmatrix}
\bold{x}_s\\
\bold{x}_d\\
\bold{y}_s\\
\bold{y}_d
\end{bmatrix}= 
\begin{bmatrix}
\bold{f}_s\\
\bold{f}_d\\
\bold{g}_s\\
\bold{g}_d
\end{bmatrix}.
\end{equation}
We group these block matrix to form a concise version,
\begin{equation}
\bold{A} = \begin{bmatrix}
\bold{A}_s   & \bold{0}    \\
\bold{0}     & \bold{A}_d  
\end{bmatrix},
\quad
\bold{B} = \begin{bmatrix}
-\bold{B}_s & \bold{0}   \\
\bold{0}   & -\bold{B}_d
\end{bmatrix},
\quad
\bold{C} = \begin{bmatrix}
\bold{C}_s &  \bold{K}   \\
\bold{K}   &  \bold{C}_d
\end{bmatrix},
\quad \bold{F} = \begin{bmatrix}
\bold{f}_s \\
\bold{f}_d
\end{bmatrix},
\quad
\bold{G} = \begin{bmatrix}
-\bold{g}_s \\
-\bold{g}_d
\end{bmatrix},
\end{equation}
\begin{equation}
\begin{bmatrix}
		  \bold{A}   & \bold{B} \\
		  \bold{B}^T & -\bold{C}
\end{bmatrix} \begin{bmatrix}
\bold{x} \\ \bold{y}
\end{bmatrix} = \begin{bmatrix}
		  \bold{F} \\
		  \bold{G}
\end{bmatrix},
\end{equation}
where $\bold{A}$ is a positive definite matrix. However,
$\bold{B}$ is usually a singular matrix when no melting region presents.

Now we take a closer look at the formed linear system.
The no melting (zero porosity $\phi=0$) region adds difficulties to solving
the formed saddle point system.
On one hand, the existence of zero porosity region guarantees that $\bold{B}$ 
matrix
formed in Darcy system does not has full column rank.
Therefore the schur complement 
\begin{equation*}
		  \mathbf{S} = -(\mathbf{C} + \mathbf{BA}^{-1}\mathbf{B}^T),
\end{equation*}
is usually not invertable directly.
On the other hand, in addition to the constant null space induced by gradient of
pressure,
the null space of $\bold{B}$ 
matrix $\text{ker}(\bold{B})$ in our darcy problem 
is not fixed due to evloution of melting (change of zero porosity region).

Taking into consideration that the solver will be implemented in parallel, 
a iterative solver is more suitable than direct solver.

\subsection{Schur Complement}
To begin with, we discuss a "direct" way to solve this derived saddle point system.
Now consider a general saddle point system~\eqref{eq:saddle} with 
symmetric positive-definite matrix $\mathbf{A}$ as
\begin{equation}
\begin{aligned}
\mathbf{A} \mathbf{x} + \mathbf{B} \mathbf{y} &= \mathbf{f},\\
\mathbf{B}^T \mathbf{x} -\mathbf{C}\mathbf{y} &= \mathbf{g}.
\end{aligned}
\end{equation}
We represent the solution vector $\mathbf{x}$ as
\begin{equation}
\mathbf{x} = \mathbf{A}^{-1} (\mathbf{f} - \mathbf{B} \mathbf{y}),
\end{equation}
where $\mathbf{A}^{-1}$ is computed with conjugate gradient algorithm.
We then eliminate $\mathbf{x}$ to represent $\mathbf{y}$ as
\begin{equation}
\mathbf{y} = \mathbf{S}^{-1}(\mathbf{g} - \mathbf{B}^T \mathbf{A}^{-1}\mathbf{f}),
\end{equation}
where $\mathbf{S}^{-1}$ is computed with Minimal Resigual algorithm~\cite{doi:10.1137/0712047} 
due to the formed Schur complement being a symmtric indefinite system.

\subsection{Uzawa Iteration}
Uzawa algorithm was originally introduced by~\cite{uzawa_arrow_hurwicz},
which gained popularity in computational fluid dynamics.
It has later been applied to the case of discretizing Stokes problem with 
Mixed Finite Element. Some inexact inner solver to accelerate uzawa iteration
was discussed by
~\cite{InexactUzawa_Elam_Golub} and ~\cite{AnalysisInexactUzawa_Bramble_Pasciak_Vassilev}.

We start by stating an abstract inexact ezawa iteration algorithm in 
Algorithm~\ref{alg:inexactuzawa}.
\begin{algorithm}
		  \caption{Inexact Uzawa Iteration\label{alg:inexactuzawa}}
		  \begin{algorithmic}[1]
					 \State Initialize $\bold{x}_0 = \bold{0}$, $\bold{y}_0 = \bold{0}$, tolence $\text{tol} = 1$ and residual $res = 0$.
					 \vspace{0.3cm}
					 \While {res $>$ tol}
					 \vspace{0.3cm}
					 \State Obtain $\tilde{\mathbf{x}} = \bold{Q}_A^{-1} (\bold{F} - (\bold{A}\bold{x}_i + \bold{B}\bold{y}_i))$
					 \vspace{0.3cm}
					 \State Update $\bold{x}_{i+1} = \bold{x}_i + \tilde{\mathbf{x}}$
					 \vspace{0.3cm}
					 \State Obtain $\tilde{\bold{y}} = \mathbf{Q}_B^{-1}(\bold{B}^T\bold{x}_{i+1} + \bold{C}\bold{y}_{i} - \bold{G})$
					 \vspace{0.3cm}
					 \State Update $\bold{y}_{i+1} = \bold{y}_i + \tilde{\bold{y}}$
					 \vspace{0.3cm}
					 \State Update $\text{res} = \norm{\tilde{\mathbf{x}}}_2 + 
					 \norm{\tilde{\mathbf{y}}}_2$
					 \vspace{0.3cm}
					 \EndWhile
		  \end{algorithmic}
\end{algorithm}
Some choose diagonal matrix $\alpha \mathbf{I}$ for precondictioner $\mathbf{Q}_A$. However, such simple matrix wouldn't generate a convergence result when applied
to our coupled problem. In the next section, we introduce a modified algorithm
that would solve our system.

\subsection{Modified Uzawa Iteration}
We will present a modified Uzawa iteration algorithm to deal with this
coupled degenerate singular system.
\begin{algorithm}
\caption{Modified Uzawa Iteration}
		  \begin{algorithmic}[1]
					 \State Initialize $\bold{x}_0 = \bold{0}$, $\bold{y}_0 = \bold{0}$, tolence $\text{tol} = 1$ and residual $res = 0$.
					 \vspace{0.3cm}
					 \While {res $>$ tol}
					 \vspace{0.3cm}
					 \State Obtain $\tilde{\mathbf{x}} = \bold{A}^{-1} (\bold{F} - (\bold{A}\bold{x}_i + \bold{B}\bold{y}_i))$
					 \vspace{0.3cm}
					 \State Update $\bold{x}_{i+1} = \bold{x}_i + \tilde{\mathbf{x}}$
					 \vspace{0.3cm}
					 \State Obtain $\bold{r} = \bold{B}^T\bold{x}_{i+1} + \bold{C}\bold{y}_{i} - \bold{G}$
					 \vspace{0.3cm}
					 \State Obtain $\tilde{\bold{y}} =\text{argmin}_{\bold{x}\in V} ||(\bold{B}^T\bold{A}^{-1}\bold{B} - \bold{C})^{-1}\tilde{\bold{y}}-\bold{r} ||$
					 \vspace{0.3cm}
					 \State Update $\bold{y}_{i+1} = \bold{y}_i + \tilde{\bold{y}}$
					 \vspace{0.3cm}

					 					 \State Update $\text{res} = \norm{\tilde{\mathbf{x}}}_2 + 
					 \norm{\tilde{\mathbf{y}}}_2$
	\vspace{0.3cm}
					 \EndWhile
\end{algorithmic}
\end{algorithm}
